% LNCS manuscript for a MICCAI workshop; eight-page main text plus references.
% Cohort-scale deformation/AQ results are inserted only when a verified
% analysis/miccai_workshop/aq_results.tex file is present.
\documentclass[runningheads]{llncs}

\usepackage[T1]{fontenc}
\usepackage{amsmath,amssymb}
\usepackage{booktabs}
\usepackage{graphicx}
\usepackage{microtype}
\usepackage{multirow}
\usepackage{tabularx}
\usepackage{xcolor}
\usepackage{url}

\graphicspath{{analysis/miccai_workshop/figures/}{analysis/arc_inpainting/}}
% Generated by analyze_arc_inpainting.py; do not edit by hand.
\newcommand{\ArcNMatched}{211}
\newcommand{\ArcNAnalyzed}{211}
\newcommand{\ArcLesionMedian}{78.6}
\newcommand{\ArcLesionQOne}{33.9}
\newcommand{\ArcLesionQThree}{154.3}
\newcommand{\ArcLeftDominant}{210}
\newcommand{\ArcRightDominant}{1}
\newcommand{\ArcUnilateralThreshold}{0.80}
\newcommand{\ArcAuditOverlapThreshold}{50}
\newcommand{\ArcSparedVolumeN}{211}
\newcommand{\ArcSparedVolumeBeforeMedian}{22.3}
\newcommand{\ArcSparedVolumeBeforeQOne}{15.9}
\newcommand{\ArcSparedVolumeBeforeQThree}{32.3}
\newcommand{\ArcSparedVolumeAfterMedian}{20.6}
\newcommand{\ArcSparedVolumeAfterQOne}{15.0}
\newcommand{\ArcSparedVolumeAfterQThree}{30.2}
\newcommand{\ArcSparedVolumeDeltaMedian}{1.1}
\newcommand{\ArcSparedVolumeDeltaCILow}{0.04}
\newcommand{\ArcSparedVolumeDeltaCIHigh}{1.74}
\newcommand{\ArcSparedVolumeImprovedN}{120}
\newcommand{\ArcSparedVolumeImprovedPct}{56.9}
\newcommand{\ArcSparedVolumeEffect}{0.17}
\newcommand{\ArcSparedVolumeP}{$0.029$}
\newcommand{\ArcSparedAreaN}{211}
\newcommand{\ArcSparedAreaBeforeMedian}{21.4}
\newcommand{\ArcSparedAreaBeforeQOne}{15.7}
\newcommand{\ArcSparedAreaBeforeQThree}{28.7}
\newcommand{\ArcSparedAreaAfterMedian}{19.0}
\newcommand{\ArcSparedAreaAfterQOne}{14.5}
\newcommand{\ArcSparedAreaAfterQThree}{26.4}
\newcommand{\ArcSparedAreaDeltaMedian}{1.2}
\newcommand{\ArcSparedAreaDeltaCILow}{0.20}
\newcommand{\ArcSparedAreaDeltaCIHigh}{1.98}
\newcommand{\ArcSparedAreaImprovedN}{121}
\newcommand{\ArcSparedAreaImprovedPct}{57.3}
\newcommand{\ArcSparedAreaEffect}{0.23}
\newcommand{\ArcSparedAreaP}{$0.008$}
\newcommand{\ArcSparedThicknessN}{211}
\newcommand{\ArcSparedThicknessBeforeMedian}{12.8}
\newcommand{\ArcSparedThicknessBeforeQOne}{10.0}
\newcommand{\ArcSparedThicknessBeforeQThree}{17.1}
\newcommand{\ArcSparedThicknessAfterMedian}{10.5}
\newcommand{\ArcSparedThicknessAfterQOne}{7.4}
\newcommand{\ArcSparedThicknessAfterQThree}{14.6}
\newcommand{\ArcSparedThicknessDeltaMedian}{2.4}
\newcommand{\ArcSparedThicknessDeltaCILow}{1.76}
\newcommand{\ArcSparedThicknessDeltaCIHigh}{2.83}
\newcommand{\ArcSparedThicknessImprovedN}{146}
\newcommand{\ArcSparedThicknessImprovedPct}{69.2}
\newcommand{\ArcSparedThicknessEffect}{0.47}
\newcommand{\ArcSparedThicknessP}{$<0.001$}
\newcommand{\ArcClinicalLeftN}{211}
\newcommand{\ArcClinicalLeftRho}{-0.20}
\newcommand{\ArcClinicalLeftCILow}{-0.34}
\newcommand{\ArcClinicalLeftCIHigh}{-0.07}
\newcommand{\ArcClinicalLeftP}{$0.006$}
\newcommand{\ArcClinicalRightN}{211}
\newcommand{\ArcClinicalRightRho}{-0.04}
\newcommand{\ArcClinicalRightCILow}{-0.15}
\newcommand{\ArcClinicalRightCIHigh}{0.06}
\newcommand{\ArcClinicalRightP}{$0.544$}
\newcommand{\ArcClinicalVolumeN}{211}
\newcommand{\ArcClinicalVolumeRho}{-0.71}
\newcommand{\ArcClinicalVolumeCILow}{-0.78}
\newcommand{\ArcClinicalVolumeCIHigh}{-0.63}
\newcommand{\ArcClinicalVolumeP}{$<0.001$}

% Completed 2026-07-18 by compare_aq_mass_effect_models.py (seed 2026).
% The additional uncertainty-vs-combined comparison uses identical subject-
% averaged absolute errors and a 5,000-sample participant bootstrap.
\newcommand{\AQAnalysisN}{210}
\newcommand{\AQClinicalMAE}{24.23}
\newcommand{\AQClinicalRMSE}{27.83}
\newcommand{\AQClinicalRSquared}{-0.020}
\newcommand{\AQClinicalCorrelation}{-0.132}
\newcommand{\AQStandardMAE}{16.42}
\newcommand{\AQStandardRMSE}{20.36}
\newcommand{\AQStandardRSquared}{0.454}
\newcommand{\AQStandardCorrelation}{0.675}
\newcommand{\AQDeformationMAE}{16.04}
\newcommand{\AQDeformationRMSE}{20.33}
\newcommand{\AQDeformationRSquared}{0.456}
\newcommand{\AQDeformationCorrelation}{0.677}
\newcommand{\AQDeformationQCMAE}{16.00}
\newcommand{\AQDeformationQCRMSE}{20.40}
\newcommand{\AQDeformationQCRSquared}{0.452}
\newcommand{\AQDeformationQCCorrelation}{0.674}
\newcommand{\AQUncertaintyMAE}{15.71}
\newcommand{\AQUncertaintyRMSE}{20.11}
\newcommand{\AQUncertaintyRSquared}{0.468}
\newcommand{\AQUncertaintyCorrelation}{0.688}
\newcommand{\AQCombinedMAE}{15.77}
\newcommand{\AQCombinedRMSE}{19.99}
\newcommand{\AQCombinedRSquared}{0.473}
\newcommand{\AQCombinedCorrelation}{0.689}
\newcommand{\AQDeformationAdvantage}{0.38}
\newcommand{\AQDeformationAdvantageCI}{$-0.50$--$1.17$}
\newcommand{\AQDeformationP}{0.023}
\newcommand{\AQDeformationQCAdvantage}{0.43}
\newcommand{\AQDeformationQCAdvantageCI}{$-0.44$--$1.26$}
\newcommand{\AQDeformationQCP}{0.023}
\newcommand{\AQUncertaintyAdvantage}{0.71}
\newcommand{\AQUncertaintyAdvantageCI}{$-0.10$--$1.49$}
\newcommand{\AQUncertaintyP}{0.008}
\newcommand{\AQCombinedAdvantage}{0.65}
\newcommand{\AQCombinedAdvantageCI}{$-0.14$--$1.43$}
\newcommand{\AQCombinedP}{0.020}
\newcommand{\AQDeformationAfterUncertaintyAdvantage}{$-0.06$}
\newcommand{\AQDeformationAfterUncertaintyCI}{$-0.78$--$0.57$}
\newcommand{\AQDeformationAfterUncertaintyP}{0.723}


\newcommand{\plesion}{p_{\mathrm{lesion}}}
\newcommand{\logJ}{\Delta\!\log J}
\newcommand{\vect}[1]{\boldsymbol{#1}}

\begin{document}

\title{From Lesion Filling to Deformation-Aware Aphasia Modeling:\
A Lesion-Aware Stroke MRI Workflow}
\titlerunning{Lesion-Aware Stroke MRI and Aphasia Modeling}

\author{Anonymized Authors}
\authorrunning{Anonymized Authors}
\institute{Anonymized Affiliations\\
\email{email@anonymized.com}}

\maketitle

\begin{abstract}
Chronic stroke lesions violate the normal-contrast and topology assumptions of
conventional structural MRI pipelines. We present an auditable workflow that
combines automatic lesion delineation, voxelwise lesion scores and entropy,
diffusion inpainting, BrainSuite morphometry, and a contralateral-normalized
lesion-associated deformation proxy. Synthetic tissue is used only as a
computational scaffold: every biological endpoint excludes the lesion and its
dilated inpainting target. In \ArcNMatched{} acquisition-matched ARC
participants, inpainting improved the stability of target-free BrainSuite
volume, cortical-area, and cortical-thickness measurements (all
Holm-adjusted tests significant). Qualitative cases spanning 7.3--216.7~mL
showed increasingly difficult reconstruction and motivated explicit visual QC.
Lesion volume was strongly associated with Western Aphasia Battery Aphasia Quotient
(partial $\rho=\ArcClinicalVolumeRho$), while left perisylvian burden retained
an association after adjustment for volume, age, and assessment delay
($\rho=\ArcClinicalLeftRho$). We further define a cross-sectional deformation
field by removing a contralaterally fitted affine component from SVReg,
subtracting mirrored contralateral residual deformation, and summarizing radial
displacement and log-Jacobian asymmetry outside the synthetic target. A
prespecified repeated nested-cross-validation analysis tests whether lesion
uncertainty and deformation improve held-out aphasia prediction beyond lesion
volume and language-network damage. Among \AQAnalysisN{} QC-passing
participants, mean repeated-CV MAE was \AQStandardMAE{} AQ points for the
conventional lesion model, \AQDeformationMAE{} with deformation,
\AQUncertaintyMAE{} with uncertainty, and \AQCombinedMAE{} with both. Every
bootstrap interval for mean advantage over the lesion model included zero.
Adding deformation after uncertainty yielded an advantage of
\AQDeformationAfterUncertaintyAdvantage{} points (95\% CI
\AQDeformationAfterUncertaintyCI; $p=\AQDeformationAfterUncertaintyP$), showing
no additive benefit. Thus, soft lesion summaries achieved the lowest MAE, but
neither extension demonstrated robust predictive superiority.

\keywords{Stroke \and lesion inpainting \and uncertainty \and deformation
\and aphasia \and structural MRI}
\end{abstract}

\section{Introduction}
Structural MRI processing systems are commonly optimized for intact anatomy.
Chronic infarcts introduce cavities, gliosis, abnormal intensities, ventricular
enlargement, and altered cortical topology. Intensity-driven normalization may
then deform the surrounding brain to explain lesion--template mismatch
\cite{Brett2001,Andersen2010}; tissue classification and cortical surfaces may
fail beyond the visible lesion. These failures matter for aphasia studies,
where large left-hemisphere lesions intersect perisylvian cortex and where
lesion volume and location are themselves strongly related to behavior.

Lesion filling replaces abnormal intensities before conventional processing.
Classical inpainting improved nonrigid registration in multiple sclerosis
\cite{Sdika2009}; modern diffusion approaches synthesize anatomically plausible
content and improve downstream segmentation of heterogeneous lesions and
cavities \cite{Ho2020,Pollak2025}. However, plausible appearance is not evidence
that invented tissue represents the participant's premorbid anatomy. Treating
generated intensities, labels, or cortical thickness as observed biology would
create an unidentifiable endpoint. Our guiding principle is therefore to use
inpainting to stabilize measurements of real tissue while excluding every
region that touches the synthetic target.

A second challenge is that a binary lesion mask discards ambiguity at the
boundary, while a third is that lesion volume alone does not describe the
geometry of the surrounding brain. Chronic lesions can continue to expand and
the ipsilesional brain can shrink years after stroke
\cite{Seghier2014,Johnson2024}. This motivates two extensions: (i) retaining the
segmentation score and its entropy, and (ii) estimating a deformation
\emph{proxy} relative to the participant's opposite hemisphere. The latter is
not physical ground-truth mass effect; with one chronic scan it may combine
collapse, ventricular expansion, remodeling, natural asymmetry, and
registration error.

This work unifies four linked studies. Study~1 evaluates whether diffusion
inpainting stabilizes BrainSuite processing outside generated anatomy and
illustrates performance across lesion burden. Study~2 tests the
construct validity of atlas-localized lesion burden against WAB Aphasia
Quotient (AQ). Study~3 produces probabilistic lesion and entropy maps. Study~4
extracts lesion-associated deformation and prespecifies a leakage-resistant AQ
prediction comparison. We report completed cohort findings for all four
studies, including repeated nested-CV comparisons on one identically filtered
cohort.

\section{Materials and Methods}

\subsection{Cohort and acquisition matching}
The ARC BIDS archive contains chronic stroke T1-weighted MRI, participant-level
clinical variables, and WAB measurements. WAB AQ summarizes spontaneous
speech, comprehension, repetition, and naming on a 0--100 scale, with lower
scores indicating more severe aphasia \cite{Kertesz1982}. The analysis used the
intersection of completed direct and inpainted BrainSuite branches: one
acquisition for each of \ArcNMatched{} participants. If multiple scans were
available, the case table retained the acquisition paired to the study's WAB
record. Median lesion volume was \ArcLesionMedian{}~mL (IQR
\ArcLesionQOne--\ArcLesionQThree); \ArcLeftDominant{} lesions were
left-dominant and \ArcRightDominant{} was right-dominant. The lesion model was
trained on ATLAS v2 data, a large expert-delineated stroke resource
\cite{Liew2022}. The present retrospective computational study used
de-identified data under the governance of the source study [details
anonymized for review].

\subsection{Lesion delineation and predictive ambiguity}
After reorientation, N4 bias correction \cite{Tustison2010}, skull stripping,
and 12-degree-of-freedom affine alignment to a 1-mm symmetric MNI template, a
3D full-resolution nnU-Net \cite{Isensee2021} generated two-class softmax
scores. We retain the lesion channel $\plesion(x)\in[0,1]$ as a float32 NIfTI
and define the hard mask by $\plesion\geq0.5$. Importantly,
$\plesion$ is a predictive score, not a hypothesis-test $p$-value. It has not
yet been calibrated on held-out manual ARC delineations; deep segmentation
models can be overconfident \cite{Guo2017,Mehrtash2020}.

Voxelwise ambiguity is the normalized Bernoulli entropy
\begin{equation}
 H(x)=-\frac{\plesion(x)\log \plesion(x)+
 [1-\plesion(x)]\log[1-\plesion(x)]}{\log 2}.
 \label{eq:entropy}
\end{equation}
It is zero near scores 0 and 1 and maximal at 0.5. This is score ambiguity, not
a decomposition of aleatoric and epistemic uncertainty. Stored summaries
include expected lesion volume $\sum_x\plesion(x)v_x$, left/right expected
volume, entropy mass $\sum_xH(x)v_x$, mean candidate and boundary entropy,
high-entropy volume, and volumes across thresholds 0.10--0.90. The original
hard lesion and the soft summaries are both retained so downstream models can
test incremental value rather than silently replacing established lesion
features.

\subsection{Diffusion inpainting and BrainSuite processing}
The hard lesion is dilated by 3~mm within the brain to form the inpainting
target. A conditional 3D U-Net denoising diffusion model, trained with irregular
synthetic masks, receives the normalized brain and target and performs 1,000
reverse steps. Generated intensities undergo robust affine and monotonic
boundary-quantile matching followed by a 3-mm physical-space feather. Voxels
outside the target are copied exactly from the registered input. One generated
sample was used per participant in the completed morphometric study.

The inpainted T1 is processed with BrainSuite v23a
\cite{Shattuck2002}: skull stripping, tissue classification, cortical surface
extraction, cortical thickness, and surface-constrained volumetric registration
(SVReg) \cite{Joshi2009}. The same native acquisition is also processed
directly. Because the deployed branches differ in preprocessing and coordinate
grid, their comparison evaluates the complete lesion-aware workflow, not the
isolated causal effect of changing target voxels.

\subsection{Study 1: downstream processing validation}
BrainSuite labels $1000+\mathrm{ROI}$ and $2000+\mathrm{ROI}$ were reduced to
canonical cortical identifiers. A homologous left/right ROI pair was eligible
only if neither side contained even one voxel of the dilated target. Thus,
generated tissue and its blending margin contributed no volume, surface area,
or thickness endpoint. Processing stability was quantified by the normalized
discordance of each eligible homologous pair,
\begin{equation}
 D(q_L,q_R)=200\frac{|q_L-q_R|}{q_L+q_R}.
 \label{eq:asymmetry}
\end{equation}
The participant-level endpoint was the median over eligible pairs. This was a
secondary internal-consistency check, not reconstruction accuracy: normal
hemispheric differences can also increase $D$. Primary workflow evidence also
included exact preservation outside the target, completion of volumetric and
surface processing, and visual review across lesion-volume percentiles.

\subsection{Study 2: aphasia construct validity}
The original binary lesion, never the generated tissue, was intersected with
post-inpainting atlas labels. A prespecified left perisylvian network contained
pars opercularis, triangularis, and orbitalis; supramarginal and angular gyri;
superior, transverse, and middle temporal gyri; and insula. Partial Spearman
correlation related left-network lesion burden to AQ while controlling for
total lesion volume, age at stroke, and $\log(1+\mathrm{WAB\ days})$. Total
volume controlled for age and assessment delay; the right homolog network was
a negative anatomical control.

\subsection{Study 3: lesion-associated deformation proxy}
SVReg's inverse coordinate map $F(\vect{x})$ gives the subject-voxel coordinate
corresponding to atlas voxel $\vect{x}$. On valid contralateral brain voxels, a
robust affine map $F_A(\vect{x})=M\vect{x}+\vect{b}$ is fitted with iterative
median-absolute-deviation outlier rejection. Removing this global component and
mapping its residual back to atlas axes gives
\begin{equation}
 \vect{u}(\vect{x})=M^{-1}\{F(\vect{x})-F_A(\vect{x})\},
 \label{eq:residual}
\end{equation}
expressed in atlas-aligned millimeters and smoothed with a 2-mm Gaussian kernel.
Let $R$ reflect position through the atlas midline and reflect a vector's
left--right component. The lesion-associated asymmetry field is
\begin{equation}
 \vect{u}_{L}(\vect{x})=\vect{u}(\vect{x})-R\vect{u}(R\vect{x}).
 \label{eq:mirror}
\end{equation}
Both members of a mirrored pair must lie in the atlas and mapped subject brain
masks. The lesion, dilated target, and reflected target are excluded. Cases
with lesion laterality below 0.80 are flagged as unsupported.

Magnitude is $m=\|\vect{u}_L\|$. Outside the lesion, the gradient of its signed
distance map gives outward unit normal $\vect{n}$; radial displacement is
$r=\vect{u}_L\cdot\vect{n}$, with $r>0$ outward and $r<0$ inward. Local volume
change uses $J=\det(I+\nabla\vect{u})$ and
$\logJ(\vect{x})=\log J(\vect{x})-\log J(R\vect{x})$, following the
deformation-based morphometry interpretation of Jacobians
\cite{Chung2001}. Nonpositive Jacobians are invalid. Magnitude, radial, and
$\logJ$ are summarized in 3--5, 5--10, 10--20, and 20--40~mm shells. Although
Jacobian deformation has been related to radiologist-rated tumor mass effect
\cite{Zacharaki2009}, a chronic infarct is not space occupying; we therefore
avoid a mechanistic pressure claim.

The direct, non-inpainted SVReg branch is normalized identically. The magnitude
$\|\vect{u}_{L,\mathrm{inp}}-\vect{u}_{L,\mathrm{raw}}\|$ is stored as
registration sensitivity, not biological deformation. This comparison and
the exclusion of synthetic voxels address the known tendency of lesions to
bias nonlinear registration \cite{Brett2001,Sdika2009}.

\subsection{Study 4: incremental AQ prediction}
Table~\ref{tab:models} defines prespecified ridge-regression models. The primary
contrast asks whether deformation improves prediction beyond age, assessment
delay, total lesion volume, laterality, and left/right language-network lesion
burden. Eight compact 3--20~mm features are used: median and 95th-percentile
magnitude; median and mean-absolute radial displacement; outward and inward
radial integrals; and positive and negative log-Jacobian integrals. A secondary
model adds raw--inpainted sensitivity and contralateral affine-fit RMSE to test
whether apparent benefit is registration-driven. If uncertainty coverage is at
least 90\%, soft volume, entropy, boundary ambiguity, and threshold-sensitivity
features are added to the conventional lesion model, followed by deformation.

\begin{table}[t]
\caption{Prespecified AQ prediction models, evaluated on one identical
QC-passing cohort. Features are cumulative from top to bottom within each
comparison.}
\label{tab:models}
\centering
\scriptsize
\begin{tabularx}{\textwidth}{p{0.28\textwidth}X}
\toprule
Model & Predictors \\
\midrule
Clinical & Age at stroke; $\log(1+\mathrm{WAB\ days})$ \\
Conventional lesion & Clinical; lesion volume; left/right language-network burden; laterality \\
Lesion + deformation & Conventional lesion; magnitude, radial, and log-Jacobian summaries \\
+ registration QC & Above; raw--inpainted field sensitivity; affine-fit RMSE \\
Lesion + uncertainty & Conventional lesion; expected volumes, entropy mass, boundary/candidate entropy, threshold volumes \\
Uncertainty + deformation & Lesion + uncertainty; deformation summaries \\
\bottomrule
\end{tabularx}
\end{table}

Eligible fields require supported laterality, folding fraction $\leq0.05$, and
at least 1,000 valid voxels 3--20~mm from the lesion. Missing predictors are
median-imputed, standardized, and fitted with ridge regression entirely inside
each training fold. Five-fold outer cross-validation is repeated 20 times; a
four-fold inner loop selects $\alpha\in10^{-4},\ldots,10^4$ by mean absolute
error (MAE). Nested evaluation avoids parameter-selection bias
\cite{Varma2006}. Outcomes are out-of-fold MAE, RMSE, $R^2$, and Pearson $r$.
The primary paired advantage is each participant's mean absolute error under
the conventional model minus that under the deformation model; positive values
favor deformation. We prespecify participant bootstrap confidence intervals,
two-sided paired Wilcoxon tests, and Holm correction across comparisons with
the conventional lesion model \cite{Holm1979}. The direct comparison of adding
deformation after uncertainty was reported separately as a prespecified
incremental contrast.

\subsection{Statistical analysis for completed endpoints}
Study~1 used two-sided paired Wilcoxon tests, participant-bootstrap 95\%
confidence intervals (10,000 resamples), and rank-biserial effect sizes.
Study~2 used bootstrap intervals for partial Spearman correlations. Holm
correction was applied separately across the three morphometric endpoints and
the three clinical associations. All scripts preserve case-level inputs,
out-of-fold predictions, exact configurations, and subject QC images.

\section{Results}

\subsection{Completed inpainting and clinical validation}
All \ArcNAnalyzed{} inpainted images were finite, shared the expected 1-mm
geometry, and were exactly unchanged outside the dilated target.
Figure~\ref{fig:outputs}A illustrates completed BrainSuite processing for
cases nearest the 10th, 50th, and 90th lesion-volume percentiles. Inpainting
provided a continuous substrate for volumetric labeling and surface extraction
at each size. The largest case also shows why synthetic tissue must not be
interpreted biologically: its filled anatomy is visibly less plausible and
requires explicit QC. The medium case was therefore used for the complete
representative-output figure below.

As a supporting target-free stability check, median tissue-volume discordance
decreased from
\ArcSparedVolumeBeforeMedian\% to \ArcSparedVolumeAfterMedian\% (median paired
improvement \ArcSparedVolumeDeltaMedian{} points, 95\% CI
\ArcSparedVolumeDeltaCILow--\ArcSparedVolumeDeltaCIHigh; rank-biserial
$r=\ArcSparedVolumeEffect$, adjusted $p$ \ArcSparedVolumeP). Area discordance
decreased from \ArcSparedAreaBeforeMedian\% to \ArcSparedAreaAfterMedian\%
(improvement \ArcSparedAreaDeltaMedian{}, CI
\ArcSparedAreaDeltaCILow--\ArcSparedAreaDeltaCIHigh; $r=\ArcSparedAreaEffect$,
adjusted $p$ \ArcSparedAreaP). Thickness discordance decreased from
\ArcSparedThicknessBeforeMedian\% to \ArcSparedThicknessAfterMedian\%
(improvement \ArcSparedThicknessDeltaMedian{}, CI
\ArcSparedThicknessDeltaCILow--\ArcSparedThicknessDeltaCIHigh;
$r=\ArcSparedThicknessEffect$, adjusted $p$ \ArcSparedThicknessP).

Total lesion volume was strongly associated with lower AQ after adjustment for
age and assessment delay (partial $\rho=\ArcClinicalVolumeRho$, 95\% CI
\ArcClinicalVolumeCILow{} to \ArcClinicalVolumeCIHigh; adjusted $p$
\ArcClinicalVolumeP). Left perisylvian burden retained an association after
additional adjustment for total volume ($\rho=\ArcClinicalLeftRho$, CI
\ArcClinicalLeftCILow{} to \ArcClinicalLeftCIHigh; adjusted $p$
\ArcClinicalLeftP), whereas the right homolog control was null
($\rho=\ArcClinicalRightRho$, adjusted $p$ \ArcClinicalRightP). The control is
limited by the cohort's near-complete left-lesion enrichment.

\subsection{Representative probability and deformation outputs}
Figure~\ref{fig:outputs}B presents the complete stored output for one
representative medium left-hemisphere case (the 50th-percentile case in
Fig.~\ref{fig:outputs}A). Expected lesion volume was
79.09~mL versus 78.73~mL at the 0.5 threshold, and entropy mass was 24.52~mL.
Entropy concentrated along the score transition, while remaining nonzero in
ambiguous internal regions. The deformation field passed prespecified QC:
lesion laterality was 1.00, folding fraction was 1.31\%, median 3--20~mm
magnitude was 5.28~mm, median radial displacement was $+0.12$~mm, and mean
absolute radial displacement was 3.15~mm. Median raw--inpainted registration
sensitivity was 3.32~mm, giving a proxy-to-sensitivity ratio of 1.59. This
figure is technical feasibility, not a subject-specific clinical interpretation.

\begin{figure}[p]
\centering
\includegraphics[height=0.49\textheight]{representative_outputs/arc_representative_brainsuite_outputs.pdf}\par
\vspace{1mm}
\includegraphics[width=0.78\textwidth]{representative_subject_outputs.pdf}
\caption{\textbf{Inpainting and derived outputs.} (A) Original/inpainted cases
nearest the 10th, 50th, and 90th lesion-volume percentiles (7.3, 78.6, and
216.7~mL), with T1, volumetric atlas labels, and midcortical labels. The large
case exposes reduced visual plausibility. (B) Complete outputs for the medium
case: observed T1/lesion, uncalibrated $\plesion$, entropy, inpainted T1,
atlas-mapped anatomy, deformation magnitude, radial displacement, and
log-Jacobian asymmetry. Generated regions are processing scaffolds; target
voxels and their mirrors are excluded from deformation summaries.}
\label{fig:outputs}
\end{figure}

\subsection{Incremental AQ prediction}
Of 211 cases matched to the clinical table, one failed the prespecified
deformation QC, leaving \AQAnalysisN{} participants; uncertainty coverage was
100\%. Table~\ref{tab:aq-results} and Fig.~\ref{fig:aq-prediction} summarize the
20 repeated nested-CV runs. The clinical-only model was poor
(MAE \AQClinicalMAE{}), whereas the conventional lesion model achieved MAE
\AQStandardMAE{}. Deformation reduced mean MAE to \AQDeformationMAE{}, and the
registration-QC-augmented sensitivity model was similar
(\AQDeformationQCMAE{}). The uncertainty model achieved the lowest MAE
(\AQUncertaintyMAE{}); the combined uncertainty--deformation model had MAE
\AQCombinedMAE{}, the lowest RMSE (\AQCombinedRMSE{}), and the highest
$R^2$ (\AQCombinedRSquared{}) and $r$ (\AQCombinedCorrelation{}), although
these differences were small.

Against the conventional lesion model, mean participant-level MAE advantages
were \AQDeformationAdvantage{} points for deformation (95\% bootstrap CI
\AQDeformationAdvantageCI; Holm-adjusted signed-rank
$p=\AQDeformationP$) and \AQUncertaintyAdvantage{} for uncertainty (CI
\AQUncertaintyAdvantageCI; $p=\AQUncertaintyP$). The combined model's advantage
was \AQCombinedAdvantage{} (CI \AQCombinedAdvantageCI;
$p=\AQCombinedP$). Although the signed-rank tests detected distributional
shifts, every bootstrap interval for the \emph{mean} advantage crossed zero.
Evidence for improvement was therefore modest rather than definitive. Most
importantly, adding deformation directly to the uncertainty model changed mean
MAE advantage by only \AQDeformationAfterUncertaintyAdvantage{} points (CI
\AQDeformationAfterUncertaintyCI; signed-rank
$p=\AQDeformationAfterUncertaintyP$), providing no evidence of additive value.

\begin{table}[t]
\caption{Repeated nested-CV AQ prediction on the identical \AQAnalysisN-subject
cohort. Values are means across 20 outer-CV repeats.}
\label{tab:aq-results}
\centering
\small
\begin{tabular}{lrrrr}
\toprule
Model & MAE & RMSE & $R^2$ & Pearson $r$ \\
\midrule
Clinical only & \AQClinicalMAE & \AQClinicalRMSE & \AQClinicalRSquared & \AQClinicalCorrelation \\
Conventional lesion & \AQStandardMAE & \AQStandardRMSE & \AQStandardRSquared & \AQStandardCorrelation \\
Lesion + deformation & \AQDeformationMAE & \AQDeformationRMSE & \AQDeformationRSquared & \AQDeformationCorrelation \\
+ deformation + registration QC & \AQDeformationQCMAE & \AQDeformationQCRMSE & \AQDeformationQCRSquared & \AQDeformationQCCorrelation \\
Lesion + uncertainty & \AQUncertaintyMAE & \AQUncertaintyRMSE & \AQUncertaintyRSquared & \AQUncertaintyCorrelation \\
Uncertainty + deformation & \AQCombinedMAE & \AQCombinedRMSE & \AQCombinedRSquared & \AQCombinedCorrelation \\
\bottomrule
\end{tabular}
\end{table}

\begin{figure}[t]
\centering
\includegraphics[width=\textwidth]{aq_prediction_comparison.pdf}
\caption{\textbf{Completed AQ model comparison.} (A) MAE distributions across
20 repeated nested-CV splits. (B) Mean participant-level MAE advantage relative
to the conventional lesion model; positive values favor the augmented model.
Bars are 95\% participant-bootstrap CIs and $p_H$ denotes the Holm-adjusted
paired Wilcoxon result. All mean-advantage intervals cross zero.}
\label{fig:aq-prediction}
\end{figure}

\section{Discussion}
The completed cohort analysis supports lesion inpainting for a narrow but useful
purpose: a continuous computational substrate that preserves observed anatomy
outside the generated region. Exact outside-target preservation and completed
processing across lesion sizes are primary; paired target-free discordance is
supporting evidence. The difficult 90th-percentile case shows that completion
alone is not success. Post-inpainting labels also localized the retained
original lesion sufficiently to reproduce an expected AQ relationship, but
neither analysis validates invented cortex.

The uncertainty representation addresses a different question. A binary mask
commits to one contour; $\plesion$ retains boundary confidence and permits
expected-volume calculations. These summaries produced the lowest mean MAE,
suggesting that contour ambiguity carries information not fully represented by
hard volume and language-network burden. The bootstrap interval nevertheless
included zero, so this is not definitive superiority. Moreover, a softmax score
is not a statistical $p$-value and its entropy is not calibrated uncertainty.
Manual ARC masks are needed for calibration, Brier score, and
threshold-conditioned Dice; ensembles would characterize model uncertainty
more directly \cite{Mehrtash2020}.

True mass effect is causal mechanical displacement from a premorbid state; our
cross-sectional proxy instead measures left--right registration geometry after
affine removal. It assumes usable hemispheric homology and is unsuitable for
bilateral lesions. Positive radial values need not mean continuing pressure,
while negative values may reflect collapse. Natural anatomy, masking, and SVReg
regularization remain confounds, so field sensitivity, folding, affine residual,
valid coverage, and subject figures are mandatory QC.

The AQ experiment is incremental, not causal. Deformation modestly lowered mean
MAE relative to the conventional lesion model, and registration-QC covariates
did not erase that pattern. However, the mean-advantage interval crossed zero,
and deformation added nothing after uncertainty. The significant signed-rank
tests and nonsignificant bootstrap mean intervals address different summaries
of a skewed subject-error distribution; together they argue against a strong
superiority claim. Including lesion volume and language-network burden prevents
deformation from receiving credit merely as a surrogate of size or side, while
nested preprocessing limits leakage.

Limitations include no premorbid/acute reference or expert anatomical ground
truth, preprocessing-grid differences between direct and inpainted branches,
natural hemispheric differences in the supporting paired endpoint, and strong left-lesion enrichment.
Only one generative sample was used, scores are uncalibrated, and repeated CV is
not external validation. Priorities are non-inpainted and normative controls,
tissue-specific GM/WM summaries, probability calibration, and external AQ
validation.

\section{Conclusion}
Lesion-aware processing enabled BrainSuite analysis across lesion sizes while
preserving clinically meaningful lesion localization; visual QC remains
essential for large lesions. Probabilistic delineation, auditable contralateral
deformation, and leakage-resistant AQ testing extend the workflow without
treating synthetic anatomy or cross-sectional deformation as ground truth. Soft
lesion summaries yielded the lowest mean prediction error; deformation showed
no additive benefit once uncertainty was present, supporting cautious use of
both as complementary rather than proven superior biomarkers.

\begin{credits}
\subsubsection{Acknowledgments.} Funding and computing resources were provided
by Anonymized.

\subsubsection{Disclosure of Interests.} None.
\end{credits}

\clearpage
\bibliographystyle{splncs04}
\bibliography{lesionaware_miccai}

\end{document}
